\documentclass[11pt]{article}
\usepackage[margin=1in]{geometry}
\usepackage[utf8]{inputenc}
\usepackage[T1]{fontenc}
\usepackage{hyperref}
\usepackage{url}
\usepackage{booktabs}
\usepackage{multirow}
\usepackage{amsfonts}
\usepackage{amsmath}
\usepackage{amssymb}
\usepackage{amsthm}
\usepackage{graphicx}
\usepackage{xcolor}
\usepackage{subcaption}
\usepackage{natbib}
\usepackage{longtable}

%% Theorem environments
\newtheorem{theorem}{Theorem}
\newtheorem{lemma}[theorem]{Lemma}
\newtheorem{corollary}[theorem]{Corollary}
\newtheorem{definition}[theorem]{Definition}
\newtheorem{proposition}[theorem]{Proposition}
\newtheorem{remark}[theorem]{Remark}

%% Macros
\newcommand{\R}{\mathbb{R}}
\newcommand{\E}{\mathbb{E}}
\newcommand{\Var}{\mathrm{Var}}
\newcommand{\SHAP}{\textsc{SHAP}}
\newcommand{\DASH}{\textsc{Dash}}
\newcommand{\GBDT}{\textsc{GBDT}}
\newcommand{\Params}{\mathbf{\Theta}}
\newcommand{\sH}{\mathcal{H}}
\newcommand{\sO}{\mathcal{O}}
\newcommand{\incomp}{\perp}

\title{Supporting Information for:\\
``The Impossibility of Faithful Explanation Under Underdetermination''}

\author{Drake Caraker, Bryan Arnold, David Rhoads}

\date{}

\begin{document}

\maketitle

\tableofcontents

\newpage

%%=======================================================================
\section{Proofs of All Propositions}
\label{sec:si-proofs}
%%=======================================================================

\subsection{Proof of Proposition 2 (Tightness)}

\begin{proof}
We construct three witnesses, one for each achievable pair.

\paragraph{(i) Faithful + Decisive (drop Stable).}
Set $E = \mathsf{explain}$. Faithfulness: for every $\theta$,
$\incomp(E(\theta), \mathsf{explain}(\theta)) =
\incomp(\mathsf{explain}(\theta), \mathsf{explain}(\theta))$, which is
false by irreflexivity of $\incomp$. Decisiveness: for every $\theta$ and
$h$, $\incomp(\mathsf{explain}(\theta), h) \Rightarrow
\incomp(E(\theta), h)$ holds by identity. Stability fails: the Rashomon
property provides $\theta_1, \theta_2$ with
$\mathsf{observe}(\theta_1) = \mathsf{observe}(\theta_2)$ and
$\incomp(\mathsf{explain}(\theta_1), \mathsf{explain}(\theta_2))$, so
$\mathsf{explain}(\theta_1) \neq \mathsf{explain}(\theta_2)$, hence
$E(\theta_1) \neq E(\theta_2)$.

\paragraph{(ii) Faithful + Stable (drop Decisive).}
Let $c \in \sH$ satisfy $\neg\,\incomp(c, \mathsf{explain}(\theta))$ for
all $\theta$ (a ``neutral'' or ``agnostic'' explanation). Set
$E(\theta) = c$ for all $\theta$. Faithfulness holds by construction.
Stability holds because $E$ is constant. Decisiveness fails because $E$
abstains on every distinction: $\incomp(\mathsf{explain}(\theta), h)$
does not imply $\incomp(c, h)$ in general.

\paragraph{(iii) Stable + Decisive (drop Faithful).}
Let $c \in \sH$ be ``maximally committal'': for every $\theta$ and $h$,
$\incomp(\mathsf{explain}(\theta), h) \Rightarrow \incomp(c, h)$. Set
$E(\theta) = c$. Stability holds (constant map). Decisiveness holds by
construction. Faithfulness fails because $c$ may satisfy
$\incomp(c, \mathsf{explain}(\theta))$ for some $\theta$---indeed, it
must, since it inherits incompatibilities from both sides of the Rashomon
witness.
\end{proof}

\subsection{Proof of Proposition 3 (Necessity)}

\begin{proof}
($\Rightarrow$) If the Rashomon property holds, Theorem~1 shows no
faithful, stable, and decisive $E$ exists.

($\Leftarrow$) If the Rashomon property fails, then for all $\theta_1,
\theta_2$ with $\mathsf{observe}(\theta_1) = \mathsf{observe}(\theta_2)$,
$\neg\,\incomp(\mathsf{explain}(\theta_1), \mathsf{explain}(\theta_2))$.
Set $E = \mathsf{explain}$.
\begin{itemize}
  \item \textbf{Faithfulness}: $\neg\,\incomp(\mathsf{explain}(\theta),
    \mathsf{explain}(\theta))$ by irreflexivity.
  \item \textbf{Decisiveness}: $\incomp(\mathsf{explain}(\theta), h)
    \Rightarrow \incomp(E(\theta), h)$ by identity.
  \item \textbf{Stability}: Suppose $\mathsf{observe}(\theta_1) =
    \mathsf{observe}(\theta_2)$. We need
    $\mathsf{explain}(\theta_1) = \mathsf{explain}(\theta_2)$. By the
    negated Rashomon property, their explanations are compatible:
    $\neg\,\incomp(\mathsf{explain}(\theta_1),
    \mathsf{explain}(\theta_2))$. Under the additional hypothesis that
    the native explanation is \emph{injective on incompatibility classes}
    (i.e., compatible explanations on the same observe-fiber are equal),
    stability follows. In the Lean formalization, this is captured by the
    \texttt{no\_rashomon\_implies\_all\_three} theorem, which assumes the
    negation of the Rashomon property directly implies equality of native
    explanations on each fiber.
\end{itemize}
\end{proof}

\subsection{Proof of Proposition 4 ($G$-Invariance Implies Stability)}

\begin{proof}
Let $\theta_1, \theta_2$ with $\mathsf{observe}(\theta_1) =
\mathsf{observe}(\theta_2)$. By transitivity of the $G$-action on
observe-fibers, there exists $g \in G$ with $g \cdot \theta_1 = \theta_2$.
By $G$-invariance, $E(\theta_2) = E(g \cdot \theta_1) = E(\theta_1)$.
\end{proof}


%%=======================================================================
\section{All Nine Instance Details}
\label{sec:si-instances}
%%=======================================================================

\subsection{Instance 1: Feature Attribution Under Collinearity}

\paragraph{Explanation system.}
Configuration space $\Params$: trained models from a fixed class
(GBDT, Lasso, neural networks), indexed by random seed.
Observable space $Y$: input--output functions.
Explanation space $\sH$: attribution vectors
$(\varphi_1, \ldots, \varphi_P) \in \R^P$.
Incompatibility: $\varphi \incomp \varphi'$ if there exist collinear
features $j, k$ with $\varphi_j > \varphi_k$ and $\varphi'_k > \varphi'_j$.

\paragraph{Rashomon property.}
Collinear features induce a Rashomon set: models that differ only in
which feature within a correlated group acts as primary
predictor~\cite{damour2022underspecification, fisher2019models}.

\paragraph{Quantitative bounds.}
GBDT: ratio $1/(1-\rho^2)$ diverges as $\rho \to 1$.
Lasso: infinite ratio (one feature zeroed).
Random forest: $O(1/\sqrt{T})$ convergence (the mechanism behind
\DASH{}).

\paragraph{Resolution.}
\DASH{}: average SHAP values over ensemble of retrained models.

\paragraph{Experiments.}
20 XGBoost models on UCI datasets, differing only in random seed.
SHAP flip rate: 33\%. Prediction equivalence: $>$99\%.

\subsection{Instance 2: Attention Maps}

\paragraph{Explanation system.}
$\Params$: neural network weight configurations.
$Y$: input--output functions.
$\sH = \Delta^T$: attention distributions over $T$ tokens.
$\incomp$: different argmax tokens.

\paragraph{Rashomon property.}
\citet{damour2022underspecification} and \citet{jainwallace2019}
demonstrate that functionally equivalent networks attend maximally to
different tokens.

\paragraph{Experiments.}
10 DistilBERT models with weight perturbations ($\sigma \in \{0.01, 0.02\}$)
on 200 SST-2 sentences. Prediction agreement: 100\%. Argmax flip rate:
60.0\% [95\% CI: 59--61\%]. Mean Kendall $\tau$: 0.30.

\paragraph{Resolution.}
Averaged attention rollout over the Rashomon set of equivalent models.

\subsection{Instance 3: Counterfactual Explanations}

\paragraph{Explanation system.}
$\Params$: model parameters defining decision boundaries.
$Y$: prediction vectors on test set.
$\sH = \mathcal{X}$: counterfactual points.
$\incomp$: opposite directions from query point
($(x'_1 - x_0)^\top(x'_2 - x_0) < 0$).

\paragraph{Rashomon property.}
Within a small loss tolerance, models can place decision boundaries on
opposite sides of the query point~\cite{semenova2022existence}.

\paragraph{Experiments.}
20 XGBoost models on German Credit. Direction flip rate: 23.5\%
[CI: 20.6--26.6\%]. Cross-model recourse validity: 61.3\%.
Most unstable feature: \texttt{checking\_status} at 43\% flip rate.

\paragraph{Resolution.}
Rashomon-robust recourse: the set of counterfactuals valid across all
models in the Rashomon set~\cite{karimi2020algorithmic}.

\subsection{Instance 4: Concept Probes (TCAV)}

\paragraph{Explanation system.}
$\Params$: weight configurations.
$Y$: input--output functions.
$\sH = \R^c$: concept activation vectors (unit directions).
$\incomp$: structurally distinct directions not related by
concept-subspace rotation.

\paragraph{Rashomon property.}
\citet{bolukbasi2016man} and \citet{ravfogel2020null} show that concept
directions vary substantially across architectures and seeds.

\paragraph{Experiments.}
15 MLP classifiers (128--64, ReLU) on \texttt{load\_digits}.
Prediction agreement: 97.2\%. Mean $|\cos|$: 0.10 (nearly orthogonal).
Concept direction instability ($1-|\cos|$): 0.90. TCAV score SD: 0.16.

\paragraph{Resolution.}
Averaged concept activation vector or subspace over the Rashomon set.

\subsection{Instance 5: Causal Discovery}

See main text, Section~3.2. The Rashomon property is derived from Markov
equivalence (Proposition~4). Resolution: CPDAG
reporting~\cite{chickering2002optimal}.

\subsection{Instance 6: Model Selection}

\paragraph{Explanation system.}
$\Params$: model classes (e.g., depth-3 tree, depth-5 tree, linear model).
$Y$: prediction accuracy on test set.
$\sH$: model identities.
$\incomp$: different model choices.

\paragraph{Rashomon property.}
The Rashomon set generically contains multiple models with near-identical
accuracy~\cite{rudin2024amazing, semenova2022existence}.

\paragraph{Experiments.}
20 XGBoost configurations on UCI datasets. Best-model flip rate: 80.0\%.
$\Delta$AUC $< 0.02$.

\paragraph{Resolution.}
Report the Rashomon set or ensemble rather than a single ``best'' model.

\subsection{Instance 7: Saliency Maps (GradCAM)}

\paragraph{Explanation system.}
$\Params$: CNN weight configurations.
$Y$: image classification function.
$\sH$: saliency heatmaps (spatial attribution maps).
$\incomp$: different peak-activation locations.

\paragraph{Rashomon property.}
\citet{adebayo2018sanity} show that saliency maps can be insensitive to
model weights, and \citet{hooker2019benchmark} demonstrate that saliency
rankings shift under model perturbation.

\paragraph{Experiments.}
10 ResNet-18 models with weight perturbations on CIFAR-10.
Peak flip rate: 9.6\%. Prediction agreement: 78.8\%.

\paragraph{Resolution.}
Averaged GradCAM heatmap over equivalent models.

\subsection{Instance 8: Token Citation Explanations}

\paragraph{Explanation system.}
$\Params$: language model weight configurations.
$Y$: next-token prediction function.
$\sH$: self-generated explanations (chain-of-thought, cited tokens).
$\incomp$: different cited evidence for the same prediction.

\paragraph{Rashomon property.}
\citet{turpin2024language} show that chain-of-thought explanations are
unfaithful: models produce different reasoning chains for the same
answer. \citet{lanham2023measuring} and \citet{ye2022unreliability}
confirm instability across model variants.

\paragraph{Experiments.}
10 DistilBERT models on SST-2. Citation flip rate: 34.5\%.
Prediction agreement: 100\%.

\paragraph{Resolution.}
Averaged attention rollout or ensemble explanation aggregation.

\subsection{Instance 9: Mechanistic Interpretability}

See main text, Section~3.3. Rashomon property: 85 valid circuits for a
single task~\cite{meloux2025everything}. Resolution: circuit equivalence
classes.


%%=======================================================================
\section{Experiment Methodology}
\label{sec:si-experiments}
%%=======================================================================

\subsection{General Protocol}

For each instance, we follow a common experimental design:
\begin{enumerate}
  \item \textbf{Train $K$ models} from independent random seeds on the
    same dataset with the same architecture and hyperparameters.
  \item \textbf{Verify functional equivalence}: confirm that all $K$
    models achieve statistically indistinguishable performance (prediction
    agreement $>$ 95\% or $\Delta$AUC $<$ 0.03).
  \item \textbf{Extract explanations}: apply the domain-specific
    explanation method to each model.
  \item \textbf{Measure instability}: compute pairwise instability
    metrics across all $\binom{K}{2}$ model pairs.
  \item \textbf{Negative control}: verify that instances without the
    Rashomon property (e.g., unique-solution linear regression) show
    0\% instability.
  \item \textbf{Bootstrap confidence intervals}: 1000 bootstrap
    resamples for all reported metrics.
\end{enumerate}

\subsection{Datasets}

\begin{itemize}
  \item \textbf{Attribution, counterfactual, model selection}: UCI German
    Credit, UCI Adult, and synthetic collinear datasets.
  \item \textbf{Attention, token citation}: Stanford Sentiment Treebank
    (SST-2), 200 sentences.
  \item \textbf{Concept probe}: \texttt{sklearn load\_digits}, 8$\times$8
    digit images.
  \item \textbf{Saliency map}: CIFAR-10, 100 test images.
\end{itemize}

\subsection{Model Architectures}

\begin{itemize}
  \item \textbf{Attribution}: XGBoost (100 trees, max depth 6).
  \item \textbf{Attention / token citation}: DistilBERT (6 layers, 12 heads,
    66M parameters) with Gaussian weight perturbation.
  \item \textbf{Counterfactual}: XGBoost on German Credit.
  \item \textbf{Concept probe}: MLP (128--64, ReLU), trained with SGD.
  \item \textbf{Model selection}: XGBoost with 20 hyperparameter
    configurations.
  \item \textbf{Saliency}: ResNet-18 with Gaussian weight perturbation.
\end{itemize}

\subsection{Perturbation Protocol for Attention and Saliency}

For transformer and CNN models, we use weight perturbation within the
zero-loss manifold rather than full retraining. Gaussian noise with
standard deviation $\sigma$ is added to all model weights. The
perturbation magnitude $\sigma$ is chosen to preserve prediction accuracy
within 1\% of the original model. This faithfully models the Rashomon
set: the perturbed models are functionally equivalent but differ in
internal representations.

\emph{Limitation}: The distribution over perturbations may differ from
the distribution induced by stochastic gradient descent. Full retraining
experiments for transformers are computationally expensive and are left
to future work.


%%=======================================================================
\section{Lean Verification Details}
\label{sec:si-lean}
%%=======================================================================

\subsection{Summary Statistics}

\begin{table}[h]
\centering
\caption{Lean~4 formalization summary.}
\label{tab:si-lean-summary}
\begin{tabular}{lrr}
\toprule
Component & Count & Notes \\
\midrule
Files & 73 & All in \texttt{UniversalImpossibility/} \\
Theorems + lemmas & 346 & Verified by \texttt{grep -c} \\
Axioms & 72 & Domain-specific; core uses 0 \\
\texttt{sorry} & 0 & Zero admitted goals \\
\midrule
Core impossibility & 1 theorem & Zero axiom dependencies \\
GBDT instance & 4 axioms & Surjectivity, first-mover, split structure \\
DASH resolution & 6 axioms & + cross-group symmetry \\
Causal instance & 3 axioms & CPDAG Markov equivalence \\
Neural net instance & 2 axioms & Overparameterization + symmetry \\
Fairness instance & 2 axioms & Unequal base rates + equal accuracy \\
\bottomrule
\end{tabular}
\end{table}

\subsection{Axiom Inventory}

The 72 axioms fall into four categories:

\paragraph{Type declarations (6 axioms).}
\texttt{Model}, \texttt{numTrees}, \texttt{numTrees\_pos},
\texttt{attribution}, \texttt{splitCount}, \texttt{firstMover}.
These declare the types and basic constants of the GBDT framework.

\paragraph{Core behavioral properties (6 axioms).}
\texttt{firstMover\_surjective}, \texttt{splitCount\_firstMover},
\texttt{splitCount\_nonFirstMover}, \texttt{proportionality\_global},
\texttt{splitCount\_crossGroup\_symmetric},
\texttt{splitCount\_crossGroup\_stable}.
These encode the structure of gradient-boosted decision trees.

\paragraph{Measure infrastructure (2 axioms).}
\texttt{modelMeasurableSpace}, \texttt{modelMeasure}. These connect the
model space to Mathlib's measure theory library.

\paragraph{Universal instance axioms (58 axioms).}
Per-instance type and Rashomon witness axioms for the nine instances
(attention, counterfactual, concept, causal, model selection, saliency,
token citation, mechanistic interpretability, fairness).

\subsection{Axiom Stratification}

The axiom dependencies are strictly stratified:
\begin{itemize}
  \item \textbf{Core impossibility} (\texttt{explanation\_impossibility}):
    \textbf{0 axioms}. The Rashomon property enters as a hypothesis, not
    an axiom.
  \item \textbf{Attribution impossibility}
    (\texttt{attribution\_impossibility}): \textbf{0 behavioral axioms}
    (only type declarations).
  \item \textbf{GBDT impossibility} (\texttt{gbdt\_impossibility\_local}):
    4 axioms (surjectivity, first-mover, non-first-mover, split-count).
  \item \textbf{Quantitative bounds} (\texttt{impossibility}):
    5 axioms (+ proportionality).
  \item \textbf{DASH resolution} (\texttt{consensus\_equity}):
    6 axioms (+ cross-group symmetry).
\end{itemize}

\subsection{Axiom Substitution Analysis}

We have verified that:
\begin{enumerate}
  \item Removing any single core behavioral axiom blocks the corresponding
    impossibility proof (each axiom is necessary).
  \item The Rashomon property can be substituted with any domain-specific
    witness (Markov equivalence, collinearity, weight-space symmetry)
    without changing the core proof.
  \item The causal instance achieves \emph{zero} axioms for the Rashomon
    property itself (it is derived from Markov equivalence as a theorem).
\end{enumerate}

\subsection{Key Theorem List}

\begin{longtable}{lll}
\toprule
Theorem name & File & Axioms \\
\midrule
\texttt{explanation\_impossibility} & ExplanationSystem.lean & 0 \\
\texttt{attribution\_impossibility} & Trilemma.lean & 0 \\
\texttt{attribution\_impossibility\_weak} & Trilemma.lean & 0 \\
\texttt{attribution\_impossibility\_abstract} & AttributionInstance.lean & 0 \\
\texttt{attention\_impossibility} & AttentionInstance.lean & 3 \\
\texttt{counterfactual\_impossibility} & CounterfactualInstance.lean & 3 \\
\texttt{concept\_impossibility} & ConceptInstance.lean & 3 \\
\texttt{causal\_instance\_impossibility} & CausalInstance.lean & 3 \\
\texttt{mech\_interp\_impossibility} & MechInterpInstance.lean & 3 \\
\texttt{gInvariant\_stable} & UniversalResolution.lean & 0 \\
\texttt{universal\_design\_space\_dichotomy} & UniversalDesignSpace.lean & 0 \\
\texttt{no\_rashomon\_all\_three} & Necessity.lean & 0 \\
\texttt{generic\_underspecification} & Ubiquity.lean & 0 \\
\texttt{consensus\_equity} & Corollary.lean & 6 \\
\texttt{gbdt\_impossibility\_local} & General.lean & 4 \\
\texttt{ratio\_tendsto\_atTop} & Ratio.lean & 5 \\
\texttt{lasso\_impossibility} & Lasso.lean & 2 \\
\texttt{causal\_rashomon\_from\_markov} & MarkovEquivalence.lean & 0 \\
\bottomrule
\end{longtable}


%%=======================================================================
\section{Regulatory Analysis}
\label{sec:si-regulatory}
%%=======================================================================

\subsection{EU AI Act Article~13 (Transparency)}

The EU AI Act~\cite{euaiact2024} requires ``appropriate transparency'' and
``meaningful information about the logic involved'' for high-risk systems.
``Meaningful'' implies faithfulness and decisiveness; ``appropriate'' across
model versions implies stability. Theorem~1 shows all three
simultaneously is impossible under the Rashomon property.

\emph{Resolution.} Report $G$-invariant explanations with explicit
uncertainty. Disclose which property is sacrificed (typically
decisiveness).

\subsection{NIST AI RMF (MAP~1.5, MEASURE~2.6)}

The NIST framework~\cite{nistai2023} requires interpretability to be
``defined, appropriate to the AI risk, and documented.'' This implicitly
assumes all three desiderata are simultaneously achievable.

\emph{Resolution.} Explicitly specify which leg of the trilemma is
prioritized for each use case and document the tradeoff in the system
card.

\subsection{OCC SR~11-7 (Model Risk Management)}

SR~11-7~\cite{occ2011sr117} requires conceptual soundness (faithfulness),
outcomes analysis across versions (stability), and effective challenge
(decisiveness). These three pillars map directly onto the trilemma.

\emph{Resolution.} Use ensemble explanations (\DASH{}) for stability and
partial faithfulness. Document decisiveness limitations. Direct effective
challenge at modeling choices rather than explanation artifacts.


%%=======================================================================
\section{Cross-Domain Connections}
\label{sec:si-cross-domain}
%%=======================================================================

\subsection{Quine--Duhem Underdetermination}

Quine~\cite{quine1951dogmas} argued that ``any statement can be held true
come what may, if we make drastic enough adjustments elsewhere in the
system.'' Duhem~\cite{duhem1906aim} held that no physical hypothesis can
be tested in isolation. Theorem~1 formalizes this: the configuration
space is the space of theories, the observable space is empirical data,
and the Rashomon property captures the existence of empirically adequate
theories with incompatible explanations.

\subsection{Economics: Partial Identification}

Manski~\cite{manski2003partial} developed partial identification theory
for settings where observational data do not pin down a unique model.
The identified set in Manski's framework corresponds to an observe-fiber
in ours; partial identification is the economic instantiation of the
Rashomon property.

\subsection{Database Theory: View Update Problem}

Bancilhon and Spyratos~\cite{bancilhon1981update} proved that translating
updates on a database view back to the base tables is inherently ambiguous
when the view is not injective. This is precisely our framework: the base
tables are configurations, the view is the observation map, and the
ambiguity is the Rashomon property.

\subsection{Medical Diagnosis: Differential Diagnosis}

When multiple conditions produce the same symptoms (the observable),
diagnosis faces the trilemma. Differential diagnosis---the standard
medical practice of listing all conditions consistent with the
symptoms---is precisely $G$-invariant resolution: it sacrifices
decisiveness (single diagnosis) for faithfulness and stability.


%%=======================================================================
\bibliography{references}
\bibliographystyle{plainnat}
%%=======================================================================

\end{document}
